\documentclass[11pt]{article}
\usepackage[letterpaper,margin=1in]{geometry}
\usepackage{amsmath,amssymb,amsthm}

\makeatletter
\def\ps@proofguide{%
  \def\@oddhead{\small\sffamily POLYNOMIAL CONTINUED FRACTIONS\hfil rm-z2-new8}%
  \def\@oddfoot{\hfil\small\thepage}%
}
\makeatother
\pagestyle{proofguide}
\setlength{\headheight}{14pt}
\setlength{\parindent}{0pt}
\setlength{\parskip}{5pt}

\newtheorem*{theorem}{Identity}
\title{A polynomial continued fraction for a transform of \(\pi^2\)}
\author{MythicalCow}
\date{}

\begin{document}
\maketitle
\thispagestyle{proofguide}
\vspace{-1.5em}

\begin{theorem}
Let \(a_n=3n^2+9n+7\) for \(n\ge 0\) and
\(b_n=-(2n^4+3n^3-3n^2-7n-3)\) for \(n\ge 1\). The polynomial continued
fraction with these coefficients (Ramanujan Machine label \texttt{rm-z2-new8})
satisfies
\[
 a_0+\cfrac{b_1}{a_1+\cfrac{b_2}{a_2+\ddots}}
 =\frac{3\pi^2+16}{16-\pi^2}.
\]
\end{theorem}

\section{Continuants and moments}
The proof evaluates the numerator and denominator after dividing both by
the same scale. A trigonometric moment controls the error in each limit.
Let \(P_n,Q_n\) be the numerator and denominator continuants. Both satisfy
\(Z_n=a_nZ_{n-1}+b_nZ_{n-2}\) for \(n\ge1\), with
\((P_{-1},P_0)=(1,7)\) and \((Q_{-1},Q_0)=(0,1)\).

Put \(L=\pi/2\), \(u(x)=1-\cos x\), and
\[
 M_n=\int_0^L(L-x)u(x)^n\,dx,\qquad
 g_n=(n+1)(2n+1),\qquad h_n=(n+1)^2.
\]
Because \(0\le u(x)\le1\) on \([0,L]\), we have
\(0\le M_n\le M_0=L^2/2=\pi^2/8\). Using
\((u')^2=2u-u^2\) and \(u''=1-u\), differentiation gives
\((u^{n+1})''=g_nu^n-h_nu^{n+1}\). The derivative of
\((L-x)(u^{n+1})'+u^{n+1}\) is \((L-x)(g_nu^n-h_nu^{n+1})\).
The expression being differentiated is \(0\) at \(x=0\) and \(1\) at
\(x=L\), so
\begin{equation}\label{eq:moment}
 g_nM_n-h_nM_{n+1}=1.
\end{equation}

\section{Matching the recurrences}
The moment relation suggests a triangular system whose solutions can be
projected onto the continuant recurrence. Consider the triangular recurrence
\(X_{n+1}=g_nX_n+Y_n\), \(Y_{n+1}=h_nY_n\), with
\((X_0,Y_0)=(s,t)\). Set
\begin{align*}
 d_n&=(n+1)(n+2)(n+5)(n^3+13n^2+22n+8),\\
 c_n&=n^4+18n^3+91n^2+162n+96.
\end{align*}
The sequence \(W_n=d_nX_n+c_nY_n\) satisfies the continuant recurrence.
Indeed, writing \(r_n=d_{n+1}+h_nc_{n+1}\), one has
\(W_{n+1}=d_{n+1}g_nX_n+r_nY_n\). Substitution of the displayed
polynomials gives
\begin{align*}
 d_{n+2}g_{n+1}g_n&=a_{n+2}d_{n+1}g_n+b_{n+2}d_n,\\
 d_{n+2}g_{n+1}+h_nr_{n+1}&=a_{n+2}r_n+b_{n+2}c_n .
\end{align*}
\newpage
Thus \(W_{n+2}=a_{n+2}W_{n+1}+b_{n+2}W_n\). At \(n=0,1\) its
values are \(80s+96t\) and \(1584s+1952t\), respectively. Matching these
initial values yields
\[
 W_n=(8s+16t)P_n+(24s-16t)Q_n .
\]

\section{The common scale and the limit}
Let \(S_n=\prod_{k=0}^{n-1}g_k\) and
\(F_n=\prod_{k=0}^{n-1}h_k\), with empty products equal to \(1\).
Equation~\eqref{eq:moment} gives
\[
 \frac{X_{n+1}}{S_{n+1}}-\frac{X_n}{S_n}
 =\frac{tF_n}{S_{n+1}}
 =t\left(\frac{F_nM_n}{S_n}-\frac{F_{n+1}M_{n+1}}{S_{n+1}}\right),
\]
where \(Y_n=tF_n\). Summing and substituting into \(W_n\) gives
\[
 \frac{W_n}{d_nS_n}=s+tM_0+
 t\frac{F_n}{S_n}\left(\frac{c_n}{d_n}-M_n\right).
\]
For \(k\ge1\), \(h_k/g_k=(k+1)/(2k+1)\le2/3\), so
\(F_n/S_n\le(2/3)^{n-1}\) for \(n\ge1\).
The factor \(c_n/d_n-M_n\) is bounded because \(c_n/d_n\to0\)
and \(0\le M_n\le M_0\).
Consequently \(W_n/(d_nS_n)\to s+tM_0\).

Taking \((s,t)=(2,3)\) and \((2,-1)\) gives \(W_n=64P_n\)
and \(W_n=64Q_n\), respectively. Since \(2-M_0>0\), the denominator
is eventually nonzero and the quotient converges to
\[
 \lim_{n\to\infty}\frac{P_n}{Q_n}
 =\frac{2+3M_0}{2-M_0}=\frac{3\pi^2+16}{16-\pi^2}.
\]
\section*{Reading the Lean proof}
In \texttt{PcfProof/PCF.lean}, \texttt{moment\_recurrence} derives
equation~\eqref{eq:moment}, and \texttt{project\_eq} matches the continuants.
The formal argument uses a reciprocal bound on \(F_n/S_n\), which
also suffices to make the error vanish. The theorem
\texttt{PCF.pcf\_target} takes the quotient of the two limits.

\end{document}
